\documentclass[11pt]{article}

\usepackage[margin=1in]{geometry}
\usepackage{amsmath,amssymb,amsthm,mathtools}
\usepackage{booktabs,longtable,array}
\usepackage{enumitem}
\usepackage{microtype}
\usepackage[hidelinks]{hyperref}
\usepackage{caption}

\newtheorem{theorem}{Theorem}[section]
\newtheorem{lemma}[theorem]{Lemma}
\newtheorem{proposition}[theorem]{Proposition}
\newtheorem{corollary}[theorem]{Corollary}
\theoremstyle{definition}
\newtheorem{definition}[theorem]{Definition}
\newtheorem{remark}[theorem]{Remark}

\newcommand{\R}{\mathbb{R}}
\newcommand{\Z}{\mathbb{Z}}
\newcommand{\D}{\mathbb{D}}
\newcommand{\vol}{\operatorname{Vol}}
\newcommand{\per}{\operatorname{Per}}
\newcommand{\supp}{\operatorname{supp}}

\title{Rigorous Bounds for Covering the Unit Disk by 100 Equal Disks}
\author{Zou Yuanqing}
\date{September 2026}

\begin{document}

\maketitle

\begin{abstract}
Let $r_D(N)$ be the smallest radius for which $N$ equal closed disks
can cover the unit disk.  We give an explicit construction of $100$
equal disks whose union contains a disk of radius $\sqrt{67}$, proving
\[
  r_D(100)\le \frac{1}{\sqrt{67}}
  =0.122169444356305223\ldots .
\]
The construction is based on a triangular lattice.  Its covering
property is reduced to a finite certificate of elementary lattice
triangles; the certificate and the verification code are included in
the repository.  We also specialize the asymptotic theorem of Birgin,
Gardenghi, and Laurain to the unit disk and collect rigorous lower and
constructive upper bounds for every $N\le 100$.  The exact global value
of $r_D(100)$ remains open in the published literature.
\end{abstract}

\tableofcontents

\section{Introduction}

The disk covering problem asks for the smallest common radius for which
$N$ congruent disks can cover the unit disk.  In the notation of this
paper,
\begin{equation}
  r_D(N)
  =
  \inf_{c_1,\dots,c_N\in\R^2}
  \max_{x\in \D}
  \min_{1\le i\le N}\|x-c_i\|_2 ,
  \qquad
  \D=\{x\in\R^2:\|x\|_2\le 1\}.
  \label{eq:definition}
\end{equation}
The exact values are known only for small $N$; the cases $N\le10$ are
classical, while the global optimality question is open in general for
larger $N$ \cite{Kershner1939,FejesToth2005,FriedmanCovering}.

The purpose of this paper is twofold.  First, we prove a verified
explicit upper bound for $N=100$.  Second, we record a general
asymptotic law and rigorous finite-$N$ bounds, so that the discussion
is not based on a black-box numerical simulation.

The main construction uses the triangular lattice generated by
\begin{equation}
  u=(\sqrt3,0),\qquad
  v=\left(\frac{\sqrt3}{2},\frac32\right).
  \label{eq:basis}
\end{equation}
For an integer pair $(i,j)$, write
\begin{equation}
  q(i,j)=i^2+ij+j^2,
  \qquad
  c_{i,j}=iu+jv.
  \label{eq:centers}
\end{equation}
We use the $97$ lattice points
\begin{equation}
  A=\{(i,j)\in\Z^2:q(i,j)\le27\}
  \label{eq:A}
\end{equation}
and add three arbitrary lattice points of the next shell to obtain
$100$ centers.  The central result is
\begin{equation}
  \D_{\sqrt{67}}\subseteq
  \bigcup_{(i,j)\in A}\overline B(c_{i,j},1).
  \label{eq:maininclusion}
\end{equation}
Thus, after scaling by $1/\sqrt{67}$, the unit disk is covered by $100$
disks of radius $1/\sqrt{67}$.

The proof of \eqref{eq:maininclusion} consists of two independent
pieces:
\begin{enumerate}[label=(\roman*)]
  \item a finite Delaunay-triangle certificate, showing that every point
        of $\D_{\sqrt{67}}$ lies in an elementary triangle having at
        least one vertex in $A$;
  \item a finite critical-ray certificate, showing that the construction
        cannot cover a disk of radius larger than $\sqrt{67}$.
\end{enumerate}
The certificates are generated by the reproducible scripts in
\texttt{scripts/}.

The paper is organized as follows.  Section~\ref{sec:prelim} establishes
the elementary scaling and area bounds.  Section~\ref{sec:construction}
proves the $N=100$ upper bound.  Section~\ref{sec:asymptotic} states the
general asymptotic theorem and specializes it to the unit disk.
Section~\ref{sec:table} describes the complete $N\le100$ table.
Section~\ref{sec:discussion} discusses the distinction between the
constructed upper bound and a global optimum.

\section{Preliminaries}
\label{sec:prelim}

Let $D_R=\{x\in\R^2:\|x\|_2\le R\}$.  We write
$\overline B(c,r)$ for the closed disk of radius $r$ centered at $c$.

\begin{lemma}[Scaling]
\label{lem:scaling}
Suppose $N$ unit disks cover $D_R$.  Then $N$ disks of radius $1/R$
cover $\D$.  Consequently,
\begin{equation}
  r_D(N)\le \frac1R .
\end{equation}
\end{lemma}

\begin{proof}
Let the unit disks have centers $c_1,\dots,c_N$ and suppose
$D_R\subseteq\bigcup_i\overline B(c_i,1)$.  Apply the linear map
$x\mapsto x/R$.  It sends $D_R$ onto $\D$ and
$\overline B(c_i,1)$ onto
$\overline B(c_i/R,1/R)$.  Hence
\[
  \D\subseteq\bigcup_i\overline B(c_i/R,1/R).
\]
This proves the claim.
\end{proof}

\begin{lemma}[Area lower bound]
\label{lem:area}
For every $N\ge1$,
\begin{equation}
  r_D(N)\ge \frac1{\sqrt N}.
\end{equation}
\end{lemma}

\begin{proof}
If $N$ disks of radius $r$ cover $\D$, then their union has area at
least that of $\D$.  Their total area is at most $N\pi r^2$, whence
$N\pi r^2\ge\pi$ and $r\ge1/\sqrt N$.
\end{proof}

\begin{remark}
For $N=100$, Lemma~\ref{lem:area} gives $r_D(100)\ge1/10$.  The
construction below gives $r_D(100)\le1/\sqrt{67}$; the exact value is
not claimed.
\end{remark}

\section{The explicit construction for \texorpdfstring{$N=100$}{N=100}}
\label{sec:construction}

\subsection{Lattice shells}

Recall the triangular lattice $L=\{iu+jv:i,j\in\Z\}$ generated by
\eqref{eq:basis}.  The quadratic form $q(i,j)=i^2+ij+j^2$ is the
squared norm of $iu+jv$ up to the factor $3$:
\begin{equation}
  \|iu+jv\|_2^2=3q(i,j).
  \label{eq:normq}
\end{equation}
The number of lattice points in each shell is as follows.

\begin{table}[ht]
\centering
\caption{Shells of the triangular lattice used in \eqref{eq:A}.}
\label{tab:shells}
\begin{tabular}{@{}c r r@{}}
\toprule
$q$ & number of lattice points & cumulative count\\
\midrule
0&1&1\\
1&6&7\\
3&6&13\\
4&6&19\\
7&12&31\\
9&6&37\\
12&6&43\\
13&12&55\\
16&6&61\\
19&12&73\\
21&12&85\\
25&6&91\\
27&6&97\\
\bottomrule
\end{tabular}
\end{table}

\begin{lemma}
\label{lem:countA}
The set $A$ defined in \eqref{eq:A} has exactly $97$ elements.
\end{lemma}

\begin{proof}
This follows from the shell count in Table~\ref{tab:shells}.  For
completeness, the shell counts can be verified directly from the
quadratic form: the numbers of representations of the listed values of
$q$ are exactly the entries in the second column.
\end{proof}

Let $T$ be any three lattice points in the next shell, for example
\begin{equation}
  T=\{(-6,2),(4,-6),(2,4)\}.
  \label{eq:T}
\end{equation}
Then $|A\cup T|=100$.  The three extra points play no role in the
covering proof; they are included only to make the number of centers
exactly $100$.

\subsection{Delaunay triangles}

The triangular lattice has a Delaunay triangulation by equilateral
triangles of side $\sqrt3$.  Every such triangle has circumradius $1$.
Therefore, if a point $x$ lies in a lattice triangle one of whose
vertices is $c\in A$, then $\|x-c\|_2\le1$.

We use the following finite certificate.

\begin{definition}
An \emph{elementary lattice triangle} is a triangle whose three
vertices are lattice points and whose three sides have length $\sqrt3$.
\end{definition}

\begin{lemma}[Finite triangle certificate]
\label{lem:trianglecert}
Every elementary lattice triangle that meets $D_{\sqrt{67}}$ has at
least one vertex in $A$.
\end{lemma}

\begin{proof}
The proof is a finite enumeration, recorded in
\texttt{data/delaunay\_triangles.csv}.  If an elementary triangle meets
$D_{\sqrt{67}}$, its diameter is $\sqrt3$, so each vertex is at
distance at most $\sqrt{67}+\sqrt3<\sqrt{67}+2$ from the origin.  Thus
only lattice points in the finite disk of radius $\sqrt{67}+2$ can
occur.

The verification script enumerates all elementary triangles with
vertices in this finite set.  It finds $204$ triangles meeting the
target disk.  Among them,
\begin{itemize}
  \item $156$ have all three vertices in $A$;
  \item $36$ have exactly two vertices in $A$;
  \item $12$ have exactly one vertex in $A$;
  \item none have all three vertices outside $A$.
\end{itemize}
The twelve boundary cases are listed in Table~\ref{tab:boundary}.
Consequently every such triangle has a vertex in $A$.
\end{proof}

\begin{table}[ht]
\centering
\caption{The twelve elementary triangles with exactly one vertex in
$A$.  The middle entry in each row is the selected $A$-vertex.}
\label{tab:boundary}
\small
\begin{tabular}{@{}lll@{}}
\toprule
triangle vertices & selected $A$-vertex & triangle vertices\\
\midrule
$(-6,1),(-6,2),(-5,1)$&$(-5,1)$&$(5,-4),(6,-5),(6,-4)$\\
$(-6,4),(-6,5),(-5,4)$&$(-5,4)$&$(4,1),(4,2),(5,1)$\\
$(-5,-1),(-4,-2),(-4,-1)$&$(-4,-1)$&$(1,-6),(1,-5),(2,-6)$\\
$(-5,6),(-4,5),(-4,6)$&$(-4,5)$&$(1,4),(1,5),(2,4)$\\
$(-2,-4),(-1,-5),(-1,-4)$&$(-1,-4)$&$(4,-6),(4,-5),(5,-6)$\\
$(-2,6),(-1,5),(-1,6)$&$(-1,5)$&$(5,-1),(6,-2),(6,-1)$\\
\bottomrule
\end{tabular}
\end{table}

\begin{theorem}[Coverage]
\label{thm:coverage}
The $97$ unit disks centered at the lattice points of $A$ cover the
disk of radius $\sqrt{67}$:
\begin{equation}
  \D_{\sqrt{67}}
  \subseteq
  \bigcup_{(i,j)\in A}\overline B(c_{i,j},1).
  \label{eq:coverage}
\end{equation}
\end{theorem}

\begin{proof}
Let $x\in\D_{\sqrt{67}}$.  Since the elementary lattice triangles tile
the plane, $x$ lies in some elementary triangle $T$.  This triangle
meets $\D_{\sqrt{67}}$ and hence, by Lemma~\ref{lem:trianglecert}, has
a vertex $c_{i,j}$ with $(i,j)\in A$.  The circumradius of an
elementary triangle is $1$, so $\|x-c_{i,j}\|_2\le1$.  Therefore
$x$ belongs to the corresponding unit disk.  This proves
\eqref{eq:coverage}.
\end{proof}

\subsection{Sharpness of the constructed radius}

Coverage alone gives only $R(100)\ge\sqrt{67}$.  We also record that the
construction does not cover a disk of any larger radius.  This is a
property of the fixed construction, not a global optimality assertion.

\begin{lemma}[Critical ray]
\label{lem:sharp}
For the construction $A$, the largest centered disk contained in the
union of the unit disks is exactly $\D_{\sqrt{67}}$.
\end{lemma}

\begin{proof}
Set
\begin{equation}
  p=\left(-\frac{7\sqrt3}{2},-\frac{11}{2}\right),
  \qquad
  e=\frac{p}{\|p\|_2},
  \qquad
  \|p\|_2=\sqrt{67}.
  \label{eq:criticalpoint}
\end{equation}
For a center $c$, the intersection of the ray $\{te:t\ge0\}$ with
$\overline B(c,1)$ is the interval
\begin{equation}
  I_c=
  \left[
  c\cdot e-\sqrt{1-\|c-(c\cdot e)e\|_2^2},
  c\cdot e+\sqrt{1-\|c-(c\cdot e)e\|_2^2}
  \right],
  \label{eq:rayinterval}
\end{equation}
whenever $\|c-(c\cdot e)e\|_2\le1$, and is empty otherwise.
The finite list of intervals for the centers in $A$ is recorded in
\texttt{data/critical\_ray\_intervals.csv}.  Their union contains
$[0,\sqrt{67}]$, but no interval on this ray extends beyond
$\sqrt{67}$.  Indeed the maximal right endpoint is attained by the
centers
\begin{equation}
  c_{2,3}=\left(\frac{7\sqrt3}{2},\frac92\right),
  \qquad
  c_{1,4}=(3\sqrt3,6),
  \label{eq:criticalcenters}
\end{equation}
and equals $\sqrt{67}$ exactly.  Hence a centered disk of radius larger
than $\sqrt{67}$ cannot be contained in the union.
\end{proof}

\begin{corollary}
\label{cor:n100}
We have
\begin{equation}
  r_D(100)\le \frac1{\sqrt{67}}
  =
  0.122169444356305223\ldots .
  \label{eq:n100upper}
\end{equation}
\end{corollary}

\begin{proof}
By Theorem~\ref{thm:coverage}, $97$ unit disks already cover
$\D_{\sqrt{67}}$.  Adding the three remaining centers can only increase
the union.  Lemma~\ref{lem:scaling} with $R=\sqrt{67}$ gives
\eqref{eq:n100upper}.
\end{proof}

\section{The general asymptotic law}
\label{sec:asymptotic}

The finite construction above is deliberately elementary.  For the
general behavior as $N$ grows, we use the following theorem of Birgin,
Gardenghi, and Laurain \cite{BirginGardenghiLaurain2024}.

\begin{theorem}[BGL asymptotic bounds]
\label{thm:bgl}
Let $A\subset\R^2$ be a bounded set satisfying the polygonal-arc
regularity class of \cite{BirginGardenghiLaurain2024}, let
$\vol(A)$ be its area, let $\per(\partial A)$ be its perimeter, and let
$\bar k$ be its number of boundary arcs.  There exists
$\bar m\in\mathbb N$ with
\[
  \bar m\ge \frac{16\pi\bar k}{3\sqrt3}
\]
such that, for every $m>\bar m$, the optimal radius $r_A^*(m)$ for
covering $A$ by $m$ minimum-radius identical disks satisfies
\begin{equation}
  r_A^*(m)
  =
  \sqrt{\frac{2\vol(A)}{3\sqrt3\,m}}
  +
  R(m),
  \qquad
  \underline R(m)\le R(m)\le \overline R(m),
  \label{eq:bgldecomp}
\end{equation}
where
\begin{align}
  \underline R(m)
  &=
  -\frac{2\per(\partial A)}{3\sqrt3\,m}
  -\frac{8\pi\bar k(2\vol(A))^{1/2}}{(3\sqrt3\,m)^{3/2}},
  \label{eq:bglower}\\
  \overline R(m)
  &=
  \frac{2\mu_m\per(\partial A)+8\pi\bar k\mu_m^2}
       {(2\vol(A)3\sqrt3\,m)^{1/2}},
  \label{eq:bgupper}\\
  \mu_m
  &=
  \frac{
    2\per(\partial A)
    +\left[
      4\per(\partial A)^2
      +2\vol(A)(3\sqrt3\,m-10\pi\bar k)
    \right]^{1/2}
  }{
    3\sqrt3\,m-16\pi\bar k
  }.
  \label{eq:mu}
\end{align}
In particular,
\begin{equation}
  r_A^*(m)
  =
  \sqrt{\frac{2\vol(A)}{3\sqrt3\,m}}
  +O\!\left(\frac1m\right).
  \label{eq:bgllaw}
\end{equation}
\end{theorem}

The theorem is stated here for reference; its proof is contained in
\cite{BirginGardenghiLaurain2024}.  We now specialize it to the unit
disk.

\begin{corollary}[Disk specialization]
\label{cor:diskbounds}
For the unit disk,
\begin{equation}
  \vol(\D)=\pi,\qquad
  \per(\partial\D)=2\pi,\qquad
  \bar k=1.
  \label{eq:diskdata}
\end{equation}
Put
\begin{equation}
  T_N=3\sqrt3\,N.
  \label{eq:TN}
\end{equation}
Then, for every $N>\bar m$,
\begin{equation}
  \underline r(N)\le r_D(N)\le \overline r(N),
  \label{eq:diskinterval}
\end{equation}
where
\begin{align}
  \underline r(N)
  &=
  \sqrt{\frac{2\pi}{T_N}}
  -\frac{4\pi}{T_N}
  -\frac{8\pi\sqrt{2\pi}}{T_N^{3/2}},
  \label{eq:disklower}\\
  \overline r(N)
  &=
  \sqrt{\frac{2\pi}{T_N}}
  +
  \frac{4\pi\mu_N+8\pi\mu_N^2}{\sqrt{2\pi T_N}},
  \label{eq:diskupper}\\
  \mu_N
  &=
  \frac{
    4\pi+\sqrt{16\pi^2+2\pi(T_N-10\pi)}
  }{
    T_N-16\pi
  }.
  \label{eq:dispmu}
\end{align}
In particular,
\begin{equation}
  r_D(N)
  =
  \sqrt{\frac{2\pi}{3\sqrt3\,N}}
  +O\!\left(\frac1N\right).
  \label{eq:disklaw}
\end{equation}
\end{corollary}

\begin{proof}
Substitute \eqref{eq:diskdata} into \eqref{eq:bglower}--\eqref{eq:mu}.
The leading term becomes
\[
  \sqrt{\frac{2\pi}{3\sqrt3\,N}},
\]
and both remainders are $O(1/N)$.  This proves the stated bounds and
\eqref{eq:disklaw}.
\end{proof}

\section{Bounds for every \texorpdfstring{$N\le100$}{N <= 100}}
\label{sec:table}

For $N\le10$, the exact values are classical and are recorded in
Table~\ref{tab:exact}.  For $N\ge11$, we use
\begin{equation}
  \max\left\{\frac1{\sqrt N},\underline r(N)\right\}
  \le
  r_D(N)
  \le
  \text{(explicit construction)}.
  \label{eq:generalbounds}
\end{equation}
The lower bound is the maximum of the area bound
(Lemma~\ref{lem:area}) and the BGL lower bound
(Corollary~\ref{cor:diskbounds}).  The upper bound is supplied by the
constructions:
\begin{itemize}
  \item the known constructions tabulated by Friedman
        \cite{FriedmanCovering} for $11\le N\le25$;
  \item the triangular-lattice shell construction described in
        Section~\ref{sec:construction} for $26\le N\le100$.
\end{itemize}
The complete machine-readable table is
\texttt{data/r\_D\_N1-100\_bounds.csv}.  Selected values are shown in
Table~\ref{tab:selected}.

\begin{table}[ht]
\centering
\caption{Exact values for $N\le10$.}
\label{tab:exact}
\begin{tabular}{@{}c l c@{}}
\toprule
$N$&value&decimal\\
\midrule
1&1&1\\
2&1&1\\
3&$\sqrt3/2$&0.866025403784438647\\
4&$\sqrt2/2$&0.707106781186547524\\
5&algebraic root&0.609382864080709655\\
6&algebraic root&0.555905211416588705\\
7&$1/2$&0.5\\
8&$1/(1+2\cos(2\pi/7))$&0.445041867912628809\\
9&$\sqrt2-1$&0.414213562373095049\\
10&$1/(1+2\cos(2\pi/9))$&0.394930843634698458\\
\bottomrule
\end{tabular}
\end{table}

\begin{table}[ht]
\centering
\caption{Selected rigorous intervals from the complete $N\le100$ table.}
\label{tab:selected}
\begin{tabular}{@{}r r r l@{}}
\toprule
$N$&lower bound&upper bound&status\\
\midrule
25&0.2&0.239146311738&known construction\\
30&0.182574185835&0.277350098113&lattice construction\\
40&0.158113883008&0.2&lattice construction\\
50&0.141421356237&0.188982236505&lattice construction\\
60&0.129099444874&0.164356535502&lattice construction\\
70&0.119522860933&0.152498570333&lattice construction\\
80&0.111803398875&0.142857142857&lattice construction\\
90&0.105409255339&0.128020416345&lattice construction\\
100&0.1&0.122169444356&lattice construction\\
\bottomrule
\end{tabular}
\end{table}

\section{Discussion and status}
\label{sec:discussion}

The construction in Section~\ref{sec:construction} proves
\[
  r_D(100)\le \frac1{\sqrt{67}}.
\]
The area bound of Lemma~\ref{lem:area} gives
\[
  r_D(100)\ge \frac1{10}.
\]
Thus every claimed exact value must lie in
\[
  [0.1,0.122169444356305223\ldots].
\]
The two endpoints do not coincide.  The exact global value of
$r_D(100)$ is not known in the published literature, and this paper
does not claim it.

The general asymptotic theorem explains the shape of the table:
the leading term is the hexagonal-lattice term, and the first
correction is of order $1/N$.  The lattice construction used here is a
finite realization of that asymptotic geometry.  For small $N$, other
configurations, notably ring configurations, can perform better; this
is why the table does not apply one formula mechanically to all
indices.

All numerical data in the paper are reproducible.  In particular, the
$N=100$ coverage proof uses a finite triangle certificate rather than a
black-box global optimization or a GPU search.

\appendix

\section{Shell count}

Formula \eqref{eq:normq} shows that the shell of index $q$ is the set
of integer pairs representing $q$ by the positive definite quadratic
form $i^2+ij+j^2$.  The entries in Table~\ref{tab:shells} are the
standard representation counts for the listed values.  The cumulative
count through $q=27$ is
\[
  1+6+6+6+12+6+6+12+6+12+12+6+6=97.
\]

\section{Finite triangle certificate}

The script \texttt{scripts/verify\_certificate.py} performs the
following finite check:
\begin{enumerate}[label=\arabic*.]
  \item generate all lattice points within distance
        $\sqrt{67}+2$ of the origin;
  \item enumerate all elementary triangles with pairwise distance
        $\sqrt3$;
  \item retain the triangles whose convex hull meets
        $D_{\sqrt{67}}$;
  \item check that each retained triangle has at least one vertex in
        $A$.
\end{enumerate}
The resulting detailed list is stored in
\texttt{data/delaunay\_triangles.csv}.  The certificate reports $204$
triangles and no triangle with all three vertices outside $A$.

\section{Critical-ray certificate}

The file \texttt{data/critical\_ray\_intervals.csv} contains all
nonempty intervals \eqref{eq:rayinterval} for the direction
\eqref{eq:criticalpoint}.  The maximal endpoint is the exact value
$\sqrt{67}$.  The equality is attained by the centers
\eqref{eq:criticalcenters}.

\section{Reproducibility}

From the repository root, run
\begin{verbatim}
python scripts/verify_certificate.py
python scripts/generate_bounds.py
powershell -ExecutionPolicy Bypass -File scripts/build-paper.ps1
\end{verbatim}
The scripts require only Python 3 and standard scientific libraries
already listed in the repository.  No GPU computation is used.

\bibliographystyle{plain}
\bibliography{references}

\end{document}
